\section{Comparison and structural interpretation}

Jo's upstream reproducible draft~\cite{Jo26} uses
\(p_j=1/3000\) and certifies \(19/5000=0.0038\), leading to
\(0.6730085279277797\ldots\).  Rademacher's subsequent reproducible
refinement~\cite{Rademacher26} keeps the same uniform pressure and strengthens
the certified local value to \(191/50000=0.00382\), leading to
\(0.6730213619501665\ldots\).  The position-weighted vector \eqref{eq:p}
has the same total pressure \(1/500\) and the existing computer-assisted
certificate proves the stronger local value \(39/10000=0.0039\).

The previous version of the present paper used the coarse replacement
\[
 P_B\le\beta\,\operatorname{span}(B)
\]
inside each block and obtained
\[
 0.6730732086087052768\ldots .
\]
The present version keeps the exact pressure coefficients \(q_r\) until the
final shifted-block average.  With no additional local computation this alone
improves the bound to
\[
 0.6731115225500757512\ldots .
\]
The adjacent-pair refinement then gives
\[
 \boxed{0.6731175265883904388\ldots}.
\]

The distinction between the two improvements is structural.  Pressure
preservation removes an avoidable global accounting loss:
\[
 \beta(m-1)\quad\longrightarrow\quad\beta(m-6).
\]
Adjacent-pair pinching supplies new information about the same block Gram
matrix without changing the local seven-point certificate.  It is precisely
this extra defect lower bound that allows the argument to pass beyond the
\(A_m<1\) regime and use \(m=450\).

The natural finite-dimensional local optimization remains
\begin{equation}\label{eq:minimax}
 \sup_{\substack{p_j\ge0\\\sum p_j=1/500}}
 \inf_{g\in\R_{\ge0}^6}
 \left(\sum_{j=1}^6p_jg_j+\mathcal E_7(g)\right).
\end{equation}
We do not claim that \eqref{eq:p} solves this minimax problem, nor that
\(m=450\) is optimal for all refinements of the present method.  The theorem
uses the already certified vector and a new analytic exploitation of the
pressure it carries.
